% KULLANILAN TEKNOLOJİLER — [H] ile sabit konum, gereksiz boşluk yok
% Preamble'da gerekli: \usepackage{float} \usepackage{placeins}

\section{KULLANILAN TEKNOLOJİLER}

Framework'ün teknoloji yığını belirlenirken kurumsal ekosistemle uyumluluk, uzun vadeli bakım kolaylığı ve bileşenler arası entegrasyon kabiliyeti ön planda tutulmuştur. Tablo~\ref{tab:teknolojiler}'de kullanılan teknolojiler ve seçim gerekçeleri özetlenmektedir.

\begin{table}[H]
\centering
\caption{Kullanılan Teknolojiler ve Seçim Gerekçeleri}
\label{tab:teknolojiler}
\small
\begin{tabular}{|l|l|p{5.8cm}|}
\hline
\textbf{Teknoloji} & \textbf{Kategori} & \textbf{Seçim Gerekçesi} \\
\hline
Java 17 & Programlama Dili & Kurumsal ekosistem uyumu, güçlü tip sistemi \\
\hline
Maven & Build Aracı & Bağımlılık yönetimi, CI/CD entegrasyonu \\
\hline
Selenium WebDriver 4 & UI Otomasyonu & Cross-browser desteği, sektör standardı \\
\hline
TestNG & Test Framework & DataProvider, suite XML, test gruplama \\
\hline
RestAssured & API Otomasyonu & HTTP doğrulama, JSON path, Java entegrasyonu \\
\hline
JDBC (H2 / SQL Server) & DB Doğrulama & In-memory ve gerçek DB senaryoları \\
\hline
ExtentReports & Raporlama & HTML rapor, screenshot, detaylı loglama \\
\hline
Jackson & Veri İşleme & JSON-Java nesne dönüşümü \\
\hline
WebDriverManager & Sürücü Yönetimi & Otomatik ChromeDriver kurulumu \\
\hline
Groq LLM & Yapay Zeka & Hızlı inference, OpenAI uyumlu API \\
\hline
Flask + scikit-learn & ML Servisi & Flaky/stabil test tahmini (opsiyonel) \\
\hline
Chart.js & Dashboard & Test geçmişi ve modül trendi görselleştirme \\
\hline
\end{tabular}
\end{table}

\subsection{Java ve Maven}
Java, kurumsal test ekosisteminin temel dili olması nedeniyle tercih edilmiştir. Selenium, TestNG ve RestAssured gibi araçlar Java ile en olgun entegrasyonu sunmaktadır. Maven ise bağımlılıkları \texttt{pom.xml} üzerinden merkezi olarak yönetmekte ve tek komutla (\texttt{mvn clean test}) tüm suite'in çalıştırılmasını sağlamaktadır.

\subsection{TestNG ve Selenium WebDriver}
TestNG, JUnit'e kıyasla veri odaklı testler için çok daha güçlü altyapı sunmaktadır. Tablo~\ref{tab:testng_junit}'de iki framework karşılaştırılmaktadır.

\begin{table}[H]
\centering
\caption{TestNG ve JUnit Karşılaştırması}
\label{tab:testng_junit}
\begin{tabular}{|l|c|c|}
\hline
\textbf{Özellik} & \textbf{TestNG} & \textbf{JUnit} \\
\hline
DataProvider & \checkmark & Sınırlı \\
\hline
Suite XML & \checkmark & \texttimes \\
\hline
Test Gruplama & \checkmark & Sınırlı \\
\hline
Paralel Çalıştırma & \checkmark & \checkmark \\
\hline
Bağımlılık Yönetimi & \checkmark & \texttimes \\
\hline
\end{tabular}
\end{table}

Selenium WebDriver cross-browser desteği, açık kaynak lisansı ve Java ekosistemiyle doğal entegrasyonu nedeniyle UI otomasyon aracı olarak seçilmiştir. \texttt{ui-regression.xml} suite dosyasında \texttt{parallel="classes"} ile sınıf düzeyinde paralel koşum desteklenmektedir.

\FloatBarrier

\subsection{RestAssured ve JDBC}
API katmanında RestAssured kütüphanesi \texttt{reqres.in} üzerinde HTTP durum kodu, yanıt süresi ve JSON path doğrulamalarını gerçekleştirmektedir. DB katmanında H2 in-memory veritabanı hızlı senaryo koşumu için; Microsoft SQL Server gerçek ortam simülasyonu ve E2E sipariş doğrulaması için kullanılmaktadır.

\subsection{Page Object Model}
Page Object Model (POM), her sayfanın element tanımlarını ve aksiyonlarını ayrı bir sınıfta kapsüller. Şekil~\ref{fig:pom}'de POM katmanlarının ilişkisi gösterilmektedir.

\begin{figure}[H]
\centering
\small
\begin{verbatim}
┌─────────────────────────────────────┐
│         TEST KATMANI                │
│   LoginUiTest.runLoginScenario()    │
└──────────────┬──────────────────────┘
               │ çağırır
┌──────────────▼──────────────────────┐
│         PAGE KATMANI                │
│   LoginPage.login(email, password)  │
└──────────────┬──────────────────────┘
               │ element bulur
┌──────────────▼──────────────────────┐
│      SELF-HEALING KATMANI           │
│   SelfHealingElementFinder.find()   │
└─────────────────────────────────────┘
\end{verbatim}
\caption{Page Object Model Katman İlişkisi}
\label{fig:pom}
\end{figure}

\subsection{Veri Odaklı Test Yaklaşımı}
Test verileri JSON dosyalarında tutulmaktadır. Bu yaklaşımın avantajları şunlardır:
\begin{itemize}
    \item Test kodu ile test verisi birbirinden tamamen ayrıştırılmıştır
    \item Yeni senaryo eklemek için yalnızca JSON dosyası güncellenir
    \item Git üzerinde veri değişiklikleri bağımsız izlenebilir
    \item Jackson kütüphanesi JSON alanlarını otomatik olarak Java nesnelerine dönüştürür
\end{itemize}

\subsection{Standartlar}
Test senaryoları iki uluslararası standart baz alınarak tasarlanmıştır:

\begin{table}[H]
\centering
\caption{Uygulanan Standartlar ve Kapsam}
\label{tab:standartlar}
\begin{tabular}{|l|l|p{4.8cm}|}
\hline
\textbf{Standart} & \textbf{Kaynak} & \textbf{Uygulama Alanı} \\
\hline
ISTQB & istqb.org & Eşdeğerlik bölümlemesi, sınır değer analizi, negatif test, durum geçiş testi \\
\hline
OWASP & owasp.org & Authentication, input validation, güvenlik senaryoları \\
\hline
\end{tabular}
\end{table}

\FloatBarrier

\subsection{Groq LLM ve Failure Analysis}
Test hatalarının kök neden analizini otomatikleştirmek için Groq LLM servisi tercih edilmiştir. \texttt{AiFailureAnalyzer} sınıfı fail anında hata mesajı, modül metadata'sı ve self-healing geçmişini Groq API'ye ileterek Türkçe kök neden analizi üretmektedir. Tablo~\ref{tab:llm}'de alternatif LLM çözümleri karşılaştırılmaktadır.

\begin{table}[H]
\centering
\caption{LLM Çözümleri Karşılaştırması}
\label{tab:llm}
\begin{tabular}{|l|c|c|c|}
\hline
\textbf{Özellik} & \textbf{Groq} & \textbf{Claude API} & \textbf{GPT-4} \\
\hline
Ücretsiz Katman & \checkmark & \texttimes & \texttimes \\
\hline
Hızlı Inference & \checkmark & \checkmark & \checkmark \\
\hline
OpenAI Uyumlu API & \checkmark & \texttimes & \texttimes \\
\hline
Kurulum Kolaylığı & Kolay & Kolay & Kolay \\
\hline
Türkçe Analiz Desteği & \checkmark & \checkmark & \checkmark \\
\hline
\end{tabular}
\end{table}

\noindent Opsiyonel olarak \texttt{ai-analysis/} dizinindeki Flask servisi, geçmiş test verilerinden eğitilmiş ML modeli ile flaky/stabil test tahminini desteklemektedir.
